2. \underline{Italian intersections}

Let $S'$ and $S$ be two smooth projective surfaces over a field $k$, and let
$f:S'\longrightarrow S$ be a birational morphism.  If $D$ and $E$ are two
divisors on $S$, then

\begin{equation*}
(D,E)=(f^*D,f^*E),
\end{equation*}

%%%%%%%%%% source scan idx 19 | folio 12 %%%%%%%%%%

as one sees, for example, by moving $D$ and $E$.

It is therefore natural to seek a ``birational'' expression for the intersection
number.  I propose to explain how the Italians proceeded.  My principal source
of information is [1], which contains a long bibliography.  See also Manin [6]
and Northcott [12], [13].

(2.0) Let $S$ be the spectrum of a regular local ring of dimension 2, let $s$ be
the closed point of $S$, and let $k$ be a subfield of $k(s)$ over which $k(s)$ is
finite.  Intersection numbers and Euler--Poincaré characteristics will be
computed over $k$.

In this section, an $S$-scheme $p:S'\longrightarrow S$ will be called a
\underline{birational transform of} $S$ if it is regular and connected, and if
$p$ is proper and induces an isomorphism
$S'\setminus p^{-1}(s)\longrightarrow S\setminus\{s\}$.  Everything said about
$S$ will also apply to the localizations of its birational transforms at a
closed point.

(2.1) Let $f:(S'',p'')\longrightarrow(S',p')$ be a morphism of $S$-schemes
between birational transforms of $S$.  It is known [5] that $f$ has a unique
factorization

\begin{equation*}
f=f_n\cdots f_1,\qquad f_i:S_{i+1}\longrightarrow S_i,
\end{equation*}
such that $S_{i+1}$ is obtained from $S_i$ by blowing up a finite set $X_i$ of
closed points, and $f_i(X_{i+1})\subset X_i$.  The set $X_i$ is the set of points
of $S_i$ at which the projection of $S''$ onto $S_i$ is not an isomorphism.

In particular, $f$ induces an isomorphism over the complement of a finite subset
of $S'$ containing the generic points of the irreducible components of
$p'^{-1}(s)$.  Hence $f^{-1}$ defines an injection

%%%%%%%%%% source scan idx 20 | folio 13 %%%%%%%%%%

from the set $\operatorname{Irr}(p'^{-1}(s))$ of irreducible components of
$p'^{-1}(s)$ into $\operatorname{Irr}(p''^{-1}(s))$.

The birational transforms of $S$ have no $S$-automorphisms.  We write $I$ for the
set of $S$-isomorphism classes of birational transforms of $S$, ordered by
domination.  The set $I$ is filtered to the right and has least element $S$.

\underline{Definition} 2.2. \underline{The set} $S^*$ \underline{of points of}
$S$ \underline{infinitely near to} $s$ \underline{is the direct-limit set}

\begin{equation*}
S^*=\varinjlim_I\operatorname{Irr}(p^{-1}(s)) \quad .
\end{equation*}
If an element $t\in S^*$ is represented on a birational transform $S'$ of $S$ by
a complete reduced irreducible curve $T$, the field
$H^0(T,\underline{O}_T)$ depends only on $t$, not on the choice of $S'$.  It is
called the \underline{residue field} $k(t)$ of $t$.

Let $S_1$ be the blow-up of $S$ along $s$.  The projective line $s^*$, the inverse
image of $\{s\}$ in $S_1$, is called the \underline{curvetta} of $s$ and defines
an element $s^*$ of $S^*$ with residue field $k(s)$.

Let $S'$ be a birational transform of $S$, let $t$ be a closed point of $S'$, and
let $S'_t$ be the spectrum of the local ring of $S'$ at $t$.  There is an evident
injection $(S'_t)^*\hookrightarrow S^*$; in particular, $t$ defines an element
$t^*$ of $S^*$ with the same residue field.  We write $(S')^*$ for the union in
$S^*$ of the $(S'_t)^*$ as $t$ ranges over the closed points of $S'$.

The following proposition follows from the description 2.1 of the birational
transforms of $S$.

%%%%%%%%%% source scan idx 21 | folio 14 %%%%%%%%%%

\underline{Proposition} 2.3. \underline{For every} $t\in S^*$,
\underline{there is a unique sequence of birational transforms}
$(S_i)_{0\leq i\leq n+1}$ \underline{of} $S$ \underline{and closed points}
$t_i\in S_i$ $(0\leq i\leq n)$ \underline{such that}

(i)\quad $S_0=S$, \underline{and} $S_{i+1}$ \underline{is the blow-up of} $S_i$
\underline{along} $\{t_i\}$;

(ii)\quad $t_{i+1}\in t_i^*$ $(0\leq i<n)$, \underline{and} $t=t_n^*$.

\underline{Moreover, for every birational transform} $(S',p')$ \underline{of}
$S$, \underline{the element} $t$ \underline{belongs to the image of}
$\operatorname{Irr}(p'^{-1}(s))$ \underline{in} $S^*$ \underline{if and only if}
$S'$ \underline{dominates} $S_{n+1}$.

The integer $n$ is called the order of $t$.  We write $S(t)$ (respectively
$S_+(t)$) for the birational transform $S_n$ (respectively $S_{n+1}$) of $S$.
By abuse of notation, we also write $t$ for the closed point $t_n$ of $S_n$, and
$t^*$ for its inverse image $t_n^*$ in $S_{n+1}$.

We equip $S^*$ with the following order:

\begin{equation*}
t_1\prec t_2
\quad\Longleftrightarrow\quad
S_+(t_2)\ \text{dominates}\ S_+(t_1) \quad .
\end{equation*}
If $t_1\prec t_2$, then $k(t_2)$ is an extension of $k(t_1)$.  The ordered set
$S^*$ has least element $s^*$, and every bounded subset of $S^*$ is finite and
totally ordered ($S^*$ is a ``tree rooted at $s^*$'').  Moreover, the order $n$ of
$t\in S^*$ is the distance in the tree $S^*$ from $s^*$ to $t$.

The following proposition makes the description 2.1 of the birational
transforms of $S$ explicit.

\underline{Proposition} 2.4. \underline{The function assigning to a birational
transform} $p:S'\to S$ \underline{the subset}
$\operatorname{Irr}(p^{-1}(s))$ \underline{of} $S^*$ \underline{induces a
bijection between} $I$ (2.1) \underline{and the finite subsets of} $S^*$
\underline{which, together with each element,}

%%%%%%%%%% source scan idx 22 | folio 15 %%%%%%%%%%

\underline{contain all its predecessors}.

Let $D$ be a divisor $\geq0$ on $S$, and let $t\in S^*$.  We say that $t$
\underline{belongs to} $D$ if the closed point $t$ of $S(t)$ (2.3) belongs to the
strict transform (the schematic closure of $D-\{s\}$) of $D$ on $S(t)$.

\underline{Definition} 2.5. \underline{Let} $t,t'\in S^*$.
\underline{We say that} $t'$ \underline{is proximate to} $t$ \underline{if}
$t'\in S_+(t)^*$ (2.2) \underline{and belongs to the curvetta} $t^*$ \underline{of}
$t$ (2.3).

If $t'$ is an immediate successor of $t$, then $t'$ is proximate to $t$.  If in a
sequence $(t_i)_{0\leq i\leq n}$ every $t_{i+1}$ is an immediate successor of
$t_i$, and if $t_n$ is proximate to $t_0$, then every $t_i$ with $0<i\leq n$ is
proximate to $t_0$, while $t_n$ is not proximate to $t_i$ for
$0<i\leq n-1$.

(2.6) Let $A$ be a Noetherian local ring of dimension $n$, let $T$ be its
spectrum, let $m$ be its maximal ideal, let $k=A/m$ be its residue field, and let
$P=\sum_{0\leq i<n}a_iX^i$ be its Hilbert--Samuel polynomial:

\begin{equation*}
P(r)=\dim_k(m^r/m^{r+1})\qquad\text{for sufficiently large }r.
\end{equation*}
Recall that the multiplicity $\mu_s$ of the closed point $s$ in $T$ is defined by

\begin{equation*}
a_{n-1}=\frac{\mu_s}{(n-1)!} \quad .
\end{equation*}
When $n=1$, this simply says

\begin{equation*}
\mu_s=\dim_k(m^r/m^{r+1})\qquad\text{for sufficiently large }r.
\end{equation*}
Geometrically, $\mu_s$ may be defined as the degree, in the projective space
$\mathbf P_k(m/m^2)$, of the special fibre of the blow-up of $T$ along the closed
subscheme $\{s\}$.

%%%%%%%%%% source scan idx 23 | folio 16 %%%%%%%%%%

Recall that if $Y$ and $Z$ are two closed subschemes of a scheme $X$, then the
strict transform of $Y$ in the blow-up $\widetilde X$ of $X$ along $Z$---the
schematic closure of $Y\setminus Z$ in $\widetilde X$---is precisely the blow-up
of $Y$ along $Y\cap Z$.

Suppose that $T$ is a closed subscheme of a regular scheme $U$.  Let
$p:U_1\longrightarrow U$ be the blow-up of $U$ along $\{s\}$, let $s^*$ be the
inverse image of $s$ in $U_1$, and let $T'$ be the strict transform of $T$ in
$U_1$.  The preceding discussion shows that the multiplicity of $T$ at $s$ is
the degree, in the projective space $s^*$, of the subscheme $s^*\cap T'$.

Suppose further that $T$ is a divisor on $U$.  Let $T^*$ be its total transform
in $U_1$ (the inverse-image divisor), and define $\nu$ by
$T^*=\nu s^*+T'$.  If $X$ is a line in $s^*$, and Euler--Poincaré
characteristics are calculated as alternating sums of dimensions over $k$, then

\begin{equation}\tag{2.6.1}
(X,s^*)=
\chi(\underline{O}_X\overset{\mathbb L}{\otimes}\underline{O}_{s^*})
=\deg_X\underline{O}(s^*)=-1,
\end{equation}
and we have seen that

\begin{equation}\tag{2.6.2}
(X,T')=
\chi(\underline{O}_X\overset{\mathbb L}{\otimes}\underline{O}_{T'})
=\deg_{s^*}(s^*\cap T')=\mu_s \quad .
\end{equation}
The projection formula

\begin{equation*}
Rp_*(\underline{O}_X\overset{\mathbb L}{\otimes}\underline{O}_{T^*})
=Rp_*(\underline{O}_X\overset{\mathbb L}{\otimes}Lp^*\underline{O}_T)
=Rp_*\underline{O}_X\overset{\mathbb L}{\otimes}\underline{O}_T
\end{equation*}
and its corollary

\begin{equation*}
(X,T^*)=\nu(X,s^*)+(X,T')=0
\end{equation*}
give, by (2.6.1) and (2.6.2), the identity

\begin{equation}\tag{2.6.3}
\mu_s=\nu \quad .
\end{equation}

%%%%%%%%%% source scan idx 24 | folio 17 %%%%%%%%%%

\underline{Proposition} 2.7. \underline{Let} $T$ \underline{be a divisor}
$\geq0$ \underline{on a regular scheme} $U$, $s$ \underline{a closed point of}
$U$, $U_1$ \underline{the blow-up of} $U$ \underline{along} $\{s\}$, $s^*$
\underline{the inverse image of} $s$ \underline{in} $U_1$, $T'$ \underline{the
strict transform of} $T$ \underline{on} $U_1$, \underline{and} $T^*$
\underline{its total transform}.

\underline{The following three quantities are equal}:

(i)\quad \underline{the multiplicity, in Samuel's sense, of the point} $s$
\underline{on} $T$;

(ii)\quad \underline{the degree, in the projective space} $s^*$,
\underline{of} $s^*\cap T'$;

(iii)\quad \underline{the multiplicity of} $s^*$ \underline{in the divisor}
$T^*$.

2.8. Return to the hypotheses of 2.0.  Let $D$ be a divisor $\geq0$ on $S$ and
let $t\in S^*$.  The \underline{multiplicity} $\mu_t(D)$ \underline{of} $D$
\underline{at} $t$ is, by definition, the multiplicity at $t$ of the strict
transform of $D$ on $S(t)$ (2.3).  The support of the function
$t\mapsto\mu_t(D)$ is the set of points infinitely near to $s$ that belong to
$D$ (2.4).  When $D$ is regular, the numbers
$\mu_t(D)[k(t):k(s)]$ are zero or one: indeed, the strict transforms $D'$ of $D$
are isomorphic to $D$, and one applies the definition (2.6) of multiplicity.

Intersection numbers can be calculated in terms of multiplicities.

\underline{Theorem} 2.9. \underline{If the divisors} $D$ \underline{and} $E$
\underline{meet only at} $s$, \underline{then}

\begin{equation*}
(D,E)=\sum_{t\in S^*}\mu_t(D)\mu_t(E)[k(t):k],
\end{equation*}
\underline{where the sum is over the points infinitely near to} $s$
\underline{which belong to both} $D$ \underline{and} $E$.
